\documentclass[quick,pro]{debate}
\motion{“磕CP”热潮折射出人们对爱情的更深信仰 / 更深怀疑}
\stance{正方}{更深信仰}
\begin{document}
\debatetitle
\begin{thesisbox}人们嗑的是最严格的爱情标准，疑的只是遇不遇得到；疑在上面，信在底下。\end{thesisbox}

\section{定义与判准（标准表述）}
\begin{fields}
\field{嗑CP}{以旁观者的身份，给虚构角色或者真实人物配对，想象并期待他们之间的浪漫关系，并从中得到情感满足。}
\field{信仰 / 怀疑}{信仰：相信真正的爱情存在、值得期待和投入。怀疑：不相信它存在，或者不相信它靠得住、值得投入。}
\field{更深}{两种心态都有，比哪一种在底下，也就是谁是另一种的来源。}
\field{判准}{哪一种心态更能解释热潮的核心特征：嗑的是什么样的爱，为什么用旁观的方式去爱；并且是另一种心态的来源。\sub{我方要证明}{嗑CP 的动力来自相信真爱存在；犹豫疑的是遇见，不是爱情本身。}\sub{对方要证明}{嗑CP 的动力来自不信真爱靠得住。}}
\field{对方提偏向定义或判准时的 30 秒口径}{对方说“信仰就得自己去谈”：那所有还没遇到对的人的单身者都不信爱情，这是拿婚恋率换掉爱情观。请回到判准的两半：嗑什么，为什么这样嗑。}
\end{fields}

\section{我方论点}
\begin{tabularx}{\linewidth}{@{}L{4.6em}Y Y L{6.4em}@{}}\toprule
\thead{标签}&\thead{一句话}&\thead{核心数据 / 学理 / 案例}&\thead{出处}\\\midrule
\textbf{嗑的是理想}&嗑的是唯一、坚守、双向、好结局；为没见过的标准投入热情，并拿它拒绝将就，这就是信&学理：浪漫信念四条（唯一、爱能找到出路、理想化、一见钟情）；研究：深访 53 名 CP 粉，CP 是“将期待具象化的载体”；案例：“双向奔赴”2023 年十大流行语；二辩新一层：58.5\% 希望互相欣赏，“宁缺毋滥”&Sprecher 和 Metts 1989；许冠文、张慧 2023；《咬文嚼字》2023；中国青年报 2024\\
\textbf{疑的是条件}&疑的是遇不遇得到，不是爱是不是真的；放下的是遇见的成本，留下的是爱情&数据：单身青年 76.8\% 想谈；没脱单原因前三：圈子单一 64.8\%、难遇三观契合 48.4\%、不擅长交流 40.7\%；案例：张钰；三辩新一层：55781 名大学生，不愿脱单比例低于不愿结婚、生育；隔着屏幕的关系是补充&复旦发展研究院 2025；中国青年报 2024；中科院心理所 2025；Horton 和 Wohl 1956\\
\bottomrule\end{tabularx}

\section{对方论点 → 一句话反驳}
\begin{tabularx}{\linewidth}{@{}Y Y Y@{}}\toprule
\thead{对方会说}&\thead{我方一句话回}&\thead{对方追问时}\\\midrule
信在远处：爱只在屏幕那头&旁观不等于怀疑，看球的人不上场也信足球&“为什么不去遇而去嗑？”→ 卡在遇见；嗑和谈不冲突\\
要的是保险：画好终点才敢投入&怕受伤是把爱当真；画好终点，是相信爱能走到终点&“当真又不去就是不信靠得住”→ 疑的是能不能遇到对的人，不是爱本身\\
许冠文说现实每一步都充满风险&同一句前半：“我们对这种良好的亲密关系赋予了很高的价值”&“决定行为的是后半”→ 后半是门口，前半是门里\\
代餐：七成多觉得嗑CP 不比恋爱差&234 份学生问卷；人只给想要的东西找代餐&“饿了只吃代餐？”→ 饭馆远，胃口还在\\
婚恋退潮：结婚登记腰斩&结婚不等于爱情；2025 年全国通办后多了 65.7 万对&“又降了呢？”→ 数字随条件起落，量不出信\\
\bottomrule\end{tabularx}

\section{必问与必答}
\begin{points}{必问（对方不答就点名、重复、不换题）}
\item 如果人们不信爱情，为什么偏偏嗑“唯一”和“双向奔赴”？（全场主问题）
\item 圈子窄，是怀疑爱情，还是怀疑遇不遇得到？
\item 是因为想要才嗑，还是因为不信才嗑？
\end{points}
\begin{points}{必答（15 秒正面回答）}
\item 信爱情为什么不去谈？ → 卡在遇见：圈子窄、时间少；嗑和谈不冲突。
\item 想象里的标准不用兑现？ → 信仰本来就不以先兑现为条件；标准被带回现实，宁缺毋滥。
\item 中青报那项调查样本多少？ → 文中未公开，我方只用它说期待的方向。
\end{points}

\section{对辩战场概要}
\begin{tabularx}{\linewidth}{@{}L{5.2em}Y Y Y@{}}\toprule
\thead{对辩}&\thead{开场问题}&\thead{目标}&\thead{收束语}\\\midrule
三辩（开第一战场）&不信爱情的人，为什么嗑专一？&内容与形式：内容说明信什么&形式说明人站在哪，内容说明人信什么\\
一辩（收拢）&先有想要，还是先有不敢？&来源：想要在底下&先有想要，才会不敢\\
二辩（结辩前）&判准的两半，对方只讲了哪一半？&守判准，不开新战场&疑在门口，信在门里\\
\bottomrule\end{tabularx}

\section{如果……就……}
\begin{contingency}
\ifthen{对方把信仰定义成“必须自己去谈”}{说出 30 秒口径，拉回判准的两半。}
\ifthen{对方拿出没预料到的数据}{先问机构、年份、问的是不是嗑CP 的人，再问测的是想不想还是信不信。}
\ifthen{论点一被打成“想象不用兑现”}{收缩到论点二：疑的对象是遇见，用判准说明底下是信。}
\end{contingency}
提醒：正二申论先回应反二；正一质询反三是全场最后一次质询，问出的东西直接进结辩；正一结辩留约 40 字临场位。质询别频繁打断，被质询先答是或否；对方回避主问题就点名、重复、不换题。

\begin{recordcard}
\record{反二申论后}{反二申论主打}{\opt{反二申论主打“想象里的标准不用兑现”}\opt{反二申论主打“想谈不等于相信”}\opt{反二申论主打不在以上两条时}}{正二申论·回应反二申论}
\record{反一立论后}{反方立论的主干}{\opt{反方立论的主干是“旁观就是怀疑”}\opt{反方立论的主干是“怕风险就是怀疑”}\opt{反方立论的主干不在以上两条时}}{正二申论·拆反方立论}
\record{一辩对辩后}{反方到此最有力的一点}{\opt{反方全场最有力的是“画好终点”那句访谈}\opt{反方全场最有力的是“想谈却不去谈”}\opt{反方最有力的一点不在以上两条时}}{正三申论·处理反方最强点}
\record{二辩对辩后}{全场主战场}{\opt{全场主战场落在“旁观算不算怀疑”}\opt{全场主战场落在“想谈却不去谈”}\opt{主战场不在以上两条时}}{正一结辩·归纳分歧}
\record{二辩对辩后}{对方全场最有力的一点}{\opt{反方全场最有力的是“画好终点”}\opt{反方全场最有力的是“怕风险就是不信靠得住”}\opt{反方最有力的一点不在以上两条时}}{正一结辩·处理最强点}
\record{反一结辩时}{反一结辩收在哪}{\opt{反一结辩收在“信在别人身上”}\opt{反一结辩收在“保险式的爱”}\opt{反一结辩的收束不在以上两条时}}{正一结辩·回应反一结辩}
\record{全场}{对方回避了哪些问题}{（写下来）}{申论与结辩里点名}
\end{recordcard}
\end{document}
